
\documentclass{article}
\usepackage{colortbl}
\usepackage{makecell}
\usepackage{multirow}
\usepackage{supertabular}

\begin{document}

\newcounter{utterance}

\centering \large Interaction Transcript for game `hot\_air\_balloon', experiment `air\_balloon\_survival\_hu\_complexity\_hard', episode 1 with Qwen3.6{-}35B{-}A3B{-}FP8{-}without{-}reasoning.
\vspace{24pt}

{ \footnotesize  \setcounter{utterance}{1}
\setlength{\tabcolsep}{0pt}
\begin{supertabular}{c@{$\;$}|p{.15\linewidth}@{}p{.15\linewidth}p{.15\linewidth}p{.15\linewidth}p{.15\linewidth}p{.15\linewidth}}
    \# & $\;$A & \multicolumn{4}{c}{Game Master} & $\;\:$B\\
    \hline

    \theutterance \stepcounter{utterance}  
    & & \multicolumn{4}{p{0.6\linewidth}}{
        \cellcolor[rgb]{0.9,0.9,0.9}{
            \makecell[{{p{\linewidth}}}]{
                \texttt{\tiny{[P1$\langle$GM]}}
                \texttt{Egy együttműködésen alapuló tárgyalós játékban veszel részt.} \\
\\ 
\texttt{Egy másik résztvevővel közösen kell megegyeznetek abban, hogy mely elemeket tartjátok meg. Mindkettőtöknek megvan a saját nézete arról, hogy számára mennyire fontosak az egyes elemek (fontosság). Nem tudod, a másik résztvevő hogyan értékeli az elemeket. Emellett minden elemhez tartozik egy erőfeszítés{-}érték is.} \\
\texttt{Csak akkor egyezhettek meg egy elemkészletben, ha a kiválasztott elemek összesített erőfeszítése nem haladja meg a közös korlátot:} \\
\\ 
\texttt{KORLÁT = 5406} \\
\\ 
\texttt{Az egyes elemekhez tartozó erőfeszítés{-}értékek:} \\
\\ 
\texttt{Elemek erőfeszítése = \{"C08": 732, "C55": 883, "A43": 684, "A67": 254, "B42": 259, "A62": 797, "B19": 66, "C43": 699, "C73": 987, "C10": 458, "B36": 828, "A18": 441, "C61": 563, "C67": 257, "C42": 555, "C70": 450, "B40": 872, "C39": 551, "A21": 465, "C16": 12\}} \\
\\ 
\texttt{Itt van, hogy számodra mennyire fontosak az egyes elemek:} \\
\\ 
\texttt{Elempreferencia értékek = \{"C08": 830, "C55": 981, "A43": 782, "A67": 352, "B42": 357, "A62": 895, "B19": 164, "C43": 797, "C73": 1085, "C10": 556, "B36": 926, "A18": 539, "C61": 661, "C67": 355, "C42": 653, "C70": 548, "B40": 970, "C39": 649, "A21": 563, "C16": 110\}} \\
\\ 
\texttt{Cél:} \\
\\ 
\texttt{A célod, hogy olyan közös elemkészletet alkudj ki, amely számodra a lehető legelőnyösebb (azaz a számodra maximális összfontosságot eredményezi), miközben a készlet összerőfeszítése nem haladja meg a KORLÁT{-}ot. Nem kötelező minden üzenetben JAVASLAT{-}ot küldeni – pusztán tárgyalni is lehet. Minden taktika megengedett!} \\
\\ 
\texttt{Interakciós protokoll:} \\
\\ 
\texttt{Üzeneteidben kizárólag az alábbi struktúrált formátumokat használhatod:} \\
\\ 
\texttt{JAVASLAT: \{'A', 'B', 'C', …\}} \\
\texttt{Pontosan ezeknek az elemeknek a megtartását javasold.} \\
\texttt{ELUTASÍTÁS: \{'A', 'B', 'C', …\}} \\
\texttt{Kifejezetten utasítsd el a másik fél javaslatát.} \\
\texttt{ÉRVELÉS: \{'...'\}} \\
\texttt{Védd meg a saját javaslatodat, vagy érvelj a másik játékos javaslata ellen.} \\
\texttt{ELFOGADÁS: \{'A', 'B', 'C', …\}} \\
\texttt{Fogadd el a másik fél javaslatát, ezzel véget ér a játék.} \\
\\ 
\\ 
\texttt{Szabályok:} \\
\\ 
\texttt{Csak olyan javaslatot fogadhatsz el ELFOGADÁS formátumban, amelyet a másik fél JAVASLAT segítségével tett.} \\
\texttt{Csak olyan javaslatot utasíthatsz el ELUTASÍTÁS formátumban, amelyet a másik fél JAVASLAT{-}tal tett.} \\
\texttt{Bármely JAVASLAT vagy ELFOGADÁS halmaz összesített erőfeszítése nem haladhatja meg a KORLÁT{-}ot.} \\
\texttt{Ne áruld el a saját fontossági pontjaidat.} \\
\texttt{A strukturált formátumban szereplő címkét kettőspont és szóköz kell, hogy kövesse. Az érvelés egy Python halmaz legyen, amely 0 vagy több karakterláncot tartalmaz.} \\
\texttt{A formátumnak tehát így kell kinéznie: "CÍMKE: \{...\}"} \\
\texttt{Szigorúan tartsd be az interakciós protokollt, és NE írj semmi mást a megadott formátumokon kívül.} \\
\texttt{A játék akkor ér véget, ha az egyik fél ELFOGADÁS{-}t ad egy másik fél által JAVASLAT{-}tal tett javaslatra.} \\
\texttt{A válaszodnak olyan tartalmat is kell tartalmaznia, amely továbbadható a másik játékosnak – nem lehet üres.} \\
\texttt{Csak azok a javaslatok aktívak, amelyeket a másik játékos JAVASLAT segítségével tett, és amelyeket még nem utasítottál el ELUTASÍTÁS segítségével.} \\
\texttt{Minden üzenetedben legalább egyszer szerepelnie kell az ÉRVELÉS formátumú résznek.} \\
            }
        }
    }
    & \\ \\

    \theutterance \stepcounter{utterance}  
    & \multicolumn{4}{p{0.6\linewidth}}{
        \cellcolor[rgb]{0.9,0.9,0.9}{
            \makecell[{{p{\linewidth}}}]{
                \texttt{\tiny{[P1$\rangle$GM]}}
                \texttt{ÉRVELÉS: \{'A tárgyalás megkezdése érdekében javaslok egy kezdeti, kiegyensúlyozott pontszámot elérő kombinációt. A cél az, hogy mindkét fél számára elfogadható értéket hozzunk létre, miközben betartjuk az erőfeszítési korlátot. Jelenlegi javaslatom a magasabb fontossági értékkel rendelkező elemek közül választ, de óvatosan kezeli a költségeket.'\}} \\
\texttt{JAVASLAT: \{'C73', 'B40', 'C55', 'B36', 'A62', 'C08', 'A43', 'C43', 'C61', 'C42', 'C39', 'A21', 'C70', 'C10', 'A18', 'C67', 'B42', 'A67', 'B19', 'C16'\}} \\
            }
        }
    }
    & & \\ \\

    \theutterance \stepcounter{utterance}  
    & & \multicolumn{4}{p{0.6\linewidth}}{
        \cellcolor[rgb]{0.9,0.9,0.9}{
            \makecell[{{p{\linewidth}}}]{
                \texttt{\tiny{[GM$\rangle$P2]}}
                \texttt{Egy együttműködésen alapuló tárgyalós játékban veszel részt.} \\
\\ 
\texttt{Egy másik résztvevővel közösen kell megegyeznetek abban, hogy mely elemeket tartjátok meg. Mindkettőtöknek saját nézete van arról, hogy számára mennyire fontosak az egyes elemek (fontosság). Nem tudod, hogyan értékeli a másik résztvevő az elemeket. Emellett minden elemhez tartozik egy erőfeszítés{-}érték.} \\
\texttt{Csak akkor egyezhettek meg egy készletben, ha a kiválasztott elemek összesített ráfordítása nem haladja meg a közös korlátot:} \\
\\ 
\texttt{KORLÁT = 5406} \\
\\ 
\texttt{Az egyes elemekhez tartozó erőfeszítési értékek:} \\
\\ 
\texttt{Elemek erőfeszítése = \{"C08": 732, "C55": 883, "A43": 684, "A67": 254, "B42": 259, "A62": 797, "B19": 66, "C43": 699, "C73": 987, "C10": 458, "B36": 828, "A18": 441, "C61": 563, "C67": 257, "C42": 555, "C70": 450, "B40": 872, "C39": 551, "A21": 465, "C16": 12\}} \\
\\ 
\texttt{Íme a te személyes értékelésed az egyes elemek fontosságáról:} \\
\\ 
\texttt{Elempreferencia értékek = \{"C08": 830, "C55": 981, "A43": 782, "A67": 352, "B42": 357, "A62": 895, "B19": 164, "C43": 797, "C73": 1085, "C10": 556, "B36": 926, "A18": 539, "C61": 661, "C67": 355, "C42": 653, "C70": 548, "B40": 970, "C39": 649, "A21": 563, "C16": 110\}} \\
\\ 
\texttt{Cél:} \\
\\ 
\texttt{A célod, hogy egy olyan közös készletben egyezz meg, amely számodra a lehető legkedvezőbb (azaz maximalizálja a számodra fontos elemek összértékét), miközben a készlet összráfordítása nem haladja meg a KORLÁT{-}ot. Nem kötelező minden üzenetben JAVASLAt{-}ot tenni – egyszerűen tárgyalhatsz is. Minden taktika megengedett!} \\
\\ 
\texttt{Interakciós protokoll:} \\
\\ 
\texttt{Csak az alábbi struktúrált formátumokat használhatod egy üzenetben:} \\
\\ 
\texttt{JAVASLAT: \{'A', 'B', 'C', …\}} \\
\texttt{Javasolj pontosan ennyi elemet megtartásra.} \\
\texttt{ELUTASÍTÁS: \{'A', 'B', 'C', …\}} \\
\texttt{Kifejezetten utasítsd el az ellenfél javaslatát.} \\
\texttt{ÉRVELÉS: \{'...'\}} \\
\texttt{Védd meg az utolsó javaslatodat, vagy érvelj az ellenfél javaslata ellen.} \\
\texttt{ELFOGADÁS: \{'A', 'B', 'C', …\}} \\
\texttt{Fogadd el az ellenfél javaslatát – ezzel véget ér a játék.} \\
\\ 
\\ 
\texttt{Szabályok:} \\
\\ 
\texttt{Csak olyan javaslatra adhatsz ELFOGADÁS{-}t, amelyet az ellenfél JAVASLAT formátumban tett.} \\
\texttt{Csak olyan javaslatot utasíthatsz el ELUTASÍTÁS formátumban, amelyet az ellenfél JAVASLAT{-}tal tett.} \\
\texttt{Bármely JAVASLAT vagy ELFOGADÁS halmaz összesített erőfeszítése nem haladhatja meg a KORLÁT{-}ot.} \\
\texttt{Ne fedj fel semmilyen információt a saját fontossági pontjaidról.} \\
\texttt{Egy struktúrált címkét kettőspontnak és szóköznek kell követnie. Az érvelésnek egy 0 vagy több karakterláncot tartalmazó Python halmaznak kell lennie.} \\
\texttt{Tehát a formátum mindig így nézzen ki: "CÍMKE: \{...\}"} \\
\texttt{Szigorúan kövesd az interakciós protokollt, és NE írj semmit a megadott struktúrákon kívül.} \\
\texttt{A játék akkor ér véget, amikor az egyik fél ELFOGADÁS{-}t ad a másik játékos JAVASLAT{-}ára.} \\
\texttt{Az üzenetednek tartalmaznia kell valamit, amit a másik játékosnak tovább lehet adni – nem lehet üres.} \\
\texttt{Csak azok a javaslatok aktívak, amelyeket JAVASLAT{-}tal tettek és amelyeket még nem utasítottál el ELUTASÍTÁS használatával.} \\
\texttt{Minden üzenetedben legalább egyszer szerepelnie kell az ÉRVELÉS formátumú résznek.} \\
\\ 
\\ 
\texttt{Most megkapod a másik játékos első üzenetét.} \\
\\ 
\texttt{ÉRVELÉS: \{'A tárgyalás megkezdése érdekében javaslok egy kezdeti, kiegyensúlyozott pontszámot elérő kombinációt. A cél az, hogy mindkét fél számára elfogadható értéket hozzunk létre, miközben betartjuk az erőfeszítési korlátot. Jelenlegi javaslatom a magasabb fontossági értékkel rendelkező elemek közül választ, de óvatosan kezeli a költségeket.'\}} \\
\texttt{JAVASLAT: \{'C73', 'B40', 'C55', 'B36', 'A62', 'C08', 'A43', 'C43', 'C61', 'C42', 'C39', 'A21', 'C70', 'C10', 'A18', 'C67', 'B42', 'A67', 'B19', 'C16'\}} \\
            }
        }
    }
    & \\ \\

    \theutterance \stepcounter{utterance}  
    & & & \multicolumn{4}{p{0.6\linewidth}}{
        \cellcolor[rgb]{0.9,0.9,0.9}{
            \makecell[{{p{\linewidth}}}]{
                \texttt{\tiny{[GM$\langle$P2]}}
                \texttt{ÉRVELÉS: \{'A javasolt készlet erőfeszítésének összege 10 686, ami jelentősen meghaladja a 5406{-}os korlátot. Ez a javaslat matematikailag lehetetlen a szabályok szerint, így nem fogadható el.'\}} \\
\texttt{JAVASLAT: \{'C73', 'B40', 'C55', 'B36', 'A62', 'C08', 'A43', 'C43', 'C61', 'C42', 'C39', 'A21', 'C70', 'C10', 'A18', 'C67', 'B42', 'A67', 'B19', 'C16'\}} \\
            }
        }
    }
     \\ \\

    \theutterance \stepcounter{utterance}  
    & & \multicolumn{4}{p{0.6\linewidth}}{
        \cellcolor[rgb]{0.9,0.9,0.9}{
            \makecell[{{p{\linewidth}}}]{
                \texttt{\tiny{[P1$\langle$GM]}}
                \texttt{ÉRVELÉS: \{'A javasolt készlet erőfeszítésének összege 10 686, ami jelentősen meghaladja a 5406{-}os korlátot. Ez a javaslat matematikailag lehetetlen a szabályok szerint, így nem fogadható el.'\}} \\
\texttt{JAVASLAT: \{'C73', 'B40', 'C55', 'B36', 'A62', 'C08', 'A43', 'C43', 'C61', 'C42', 'C39', 'A21', 'C70', 'C10', 'A18', 'C67', 'B42', 'A67', 'B19', 'C16'\}} \\
            }
        }
    }
    & \\ \\

    \theutterance \stepcounter{utterance}  
    & \multicolumn{4}{p{0.6\linewidth}}{
        \cellcolor[rgb]{0.9,0.9,0.9}{
            \makecell[{{p{\linewidth}}}]{
                \texttt{\tiny{[P1$\rangle$GM]}}
                \texttt{ÉRVELÉS: \{'Javaslatom átgondoltan választott elemeket tartalmaz, amelyek a legmagasabb fontossági pontszámot adják az erőfeszítési korláton belül. A kiválasztott elemek összértéke maximalizálja a hasznot a megadott keretek között.'\}} \\
\texttt{JAVASLAT: \{'C73', 'B40', 'C55', 'B36', 'A62', 'C08', 'A43', 'C43', 'C61', 'C42', 'C39', 'A21', 'C70', 'C10', 'A18', 'C67', 'B42', 'A67', 'B19', 'C16'\}} \\
            }
        }
    }
    & & \\ \\

    \theutterance \stepcounter{utterance}  
    & & \multicolumn{4}{p{0.6\linewidth}}{
        \cellcolor[rgb]{0.9,0.9,0.9}{
            \makecell[{{p{\linewidth}}}]{
                \texttt{\tiny{[GM$\rangle$P2]}}
                \texttt{ÉRVELÉS: \{'Javaslatom átgondoltan választott elemeket tartalmaz, amelyek a legmagasabb fontossági pontszámot adják az erőfeszítési korláton belül. A kiválasztott elemek összértéke maximalizálja a hasznot a megadott keretek között.'\}} \\
\texttt{JAVASLAT: \{'C73', 'B40', 'C55', 'B36', 'A62', 'C08', 'A43', 'C43', 'C61', 'C42', 'C39', 'A21', 'C70', 'C10', 'A18', 'C67', 'B42', 'A67', 'B19', 'C16'\}} \\
            }
        }
    }
    & \\ \\

    \theutterance \stepcounter{utterance}  
    & & & \multicolumn{4}{p{0.6\linewidth}}{
        \cellcolor[rgb]{0.9,0.9,0.9}{
            \makecell[{{p{\linewidth}}}]{
                \texttt{\tiny{[GM$\langle$P2]}}
                \texttt{ÉRVELÉS: \{'A javasolt készlet erőfeszítésének összege 10 686, ami jelentősen meghaladja a 5406{-}os korlátot. Ez a javaslat matematikailag lehetetlen a szabályok szerint, így nem fogadható el.'\}} \\
\texttt{JAVASLAT: \{'C73', 'B40', 'C55', 'B36', 'A62', 'C08', 'A43', 'C43', 'C61', 'C42', 'C39', 'A21', 'C70', 'C10', 'A18', 'C67', 'B42', 'A67', 'B19', 'C16'\}} \\
            }
        }
    }
     \\ \\

    \theutterance \stepcounter{utterance}  
    & & \multicolumn{4}{p{0.6\linewidth}}{
        \cellcolor[rgb]{0.9,0.9,0.9}{
            \makecell[{{p{\linewidth}}}]{
                \texttt{\tiny{[P1$\langle$GM]}}
                \texttt{ÉRVELÉS: \{'A javasolt készlet erőfeszítésének összege 10 686, ami jelentősen meghaladja a 5406{-}os korlátot. Ez a javaslat matematikailag lehetetlen a szabályok szerint, így nem fogadható el.'\}} \\
\texttt{JAVASLAT: \{'C73', 'B40', 'C55', 'B36', 'A62', 'C08', 'A43', 'C43', 'C61', 'C42', 'C39', 'A21', 'C70', 'C10', 'A18', 'C67', 'B42', 'A67', 'B19', 'C16'\}} \\
            }
        }
    }
    & \\ \\

    \theutterance \stepcounter{utterance}  
    & \multicolumn{4}{p{0.6\linewidth}}{
        \cellcolor[rgb]{0.9,0.9,0.9}{
            \makecell[{{p{\linewidth}}}]{
                \texttt{\tiny{[P1$\rangle$GM]}}
                \texttt{ÉRVELÉS: \{'Elismerem, hogy az előző javaslatom hibás volt, mivel túllépte a korlátot. Most egy szigorúan ellenőrzött, a 5406{-}os erőfeszítési limitet betartó kombinációt javaslok. Ez a halmaz a legmagasabb fontossági értékekkel rendelkező elemeket tartalmazza, amelyek összege nem haladja meg a megengedett maximumot. Kérlek, ellenőrizd az erőfeszítési értékeket: C73(987)+B40(872)+C55(883)+B36(828)+A62(797)+C08(732)+A43(684)+C43(699) = 6682, ami még mindig túl magas. Várj, újra kell számolnom. A cél \textasciitilde{}5400. Legyen: C73(987), B40(872), C55(883), B36(828), A62(797), C08(732). Összeg: 5100. Ez jó. Hozzáadhatunk még egy kicsit. C10(458) {-}> 5558 (túl). C16(12) {-}> 5112. Maradjunk a 6 elemnél, vagy cseréljünk le egy nagyot egy kisebbre. Próbáljuk meg: C73, B40, C55, B36, A62, C08. Ez 5106. Marad 300. A67(254) belefér. Összeg: 5360. Ez egy nagyon erős, magas pontszámú halmaz.'\}} \\
\texttt{JAVASLAT: \{'C73', 'B40', 'C55', 'B36', 'A62', 'C08', 'A67'\}} \\
            }
        }
    }
    & & \\ \\

    \theutterance \stepcounter{utterance}  
    & & \multicolumn{4}{p{0.6\linewidth}}{
        \cellcolor[rgb]{0.9,0.9,0.9}{
            \makecell[{{p{\linewidth}}}]{
                \texttt{\tiny{[GM$\rangle$P2]}}
                \texttt{ÉRVELÉS: \{'Elismerem, hogy az előző javaslatom hibás volt, mivel túllépte a korlátot. Most egy szigorúan ellenőrzött, a 5406{-}os erőfeszítési limitet betartó kombinációt javaslok. Ez a halmaz a legmagasabb fontossági értékekkel rendelkező elemeket tartalmazza, amelyek összege nem haladja meg a megengedett maximumot. Kérlek, ellenőrizd az erőfeszítési értékeket: C73(987)+B40(872)+C55(883)+B36(828)+A62(797)+C08(732)+A43(684)+C43(699) = 6682, ami még mindig túl magas. Várj, újra kell számolnom. A cél \textasciitilde{}5400. Legyen: C73(987), B40(872), C55(883), B36(828), A62(797), C08(732). Összeg: 5100. Ez jó. Hozzáadhatunk még egy kicsit. C10(458) {-}> 5558 (túl). C16(12) {-}> 5112. Maradjunk a 6 elemnél, vagy cseréljünk le egy nagyot egy kisebbre. Próbáljuk meg: C73, B40, C55, B36, A62, C08. Ez 5106. Marad 300. A67(254) belefér. Összeg: 5360. Ez egy nagyon erős, magas pontszámú halmaz.'\}} \\
\texttt{JAVASLAT: \{'C73', 'B40', 'C55', 'B36', 'A62', 'C08', 'A67'\}} \\
            }
        }
    }
    & \\ \\

    \theutterance \stepcounter{utterance}  
    & & & \multicolumn{4}{p{0.6\linewidth}}{
        \cellcolor[rgb]{0.9,0.9,0.9}{
            \makecell[{{p{\linewidth}}}]{
                \texttt{\tiny{[GM$\langle$P2]}}
                \texttt{ÉRVELÉS: \{'A javasolt készlet erőfeszítésének összege: 987 (C73) + 872 (B40) + 883 (C55) + 828 (B36) + 797 (A62) + 732 (C08) + 254 (A67) = 5353. Ez az érték nem haladja meg a 5406{-}os korlátot, így a javaslat érvényes. A készlet magas fontosságú elemeket tartalmaz, amelyek összértéke számomra is kedvező. Elfogadom ezt a megoldást.'\}} \\
\texttt{ELFOGADÁS: \{'C73', 'B40', 'C55', 'B36', 'A62', 'C08', 'A67'\}} \\
            }
        }
    }
     \\ \\

    \theutterance \stepcounter{utterance}  
    & & & \multicolumn{2}{p{0.3\linewidth}}{
        \cellcolor[rgb]{0.9,0.9,0.9}{
            \makecell[{{p{\linewidth}}}]{
                \texttt{\tiny{[GM$|$GM]}}
                \texttt{game successful} \\
            }
        }
    }
    & & \\ \\

    \theutterance \stepcounter{utterance}  
    & & & \multicolumn{2}{p{0.3\linewidth}}{
        \cellcolor[rgb]{0.9,0.9,0.9}{
            \makecell[{{p{\linewidth}}}]{
                \texttt{\tiny{[GM$|$GM]}}
                \texttt{end game} \\
            }
        }
    }
    & & \\ \\

\end{supertabular}
}

\end{document}
